At precisely these quartet degrees DMR30--35 supplies the
complete evaluated bound
\[
 \frac{\Xi_{k,N_a}}{kq_k}
 =O_h\!\left[
 \left(\frac{\log(k+2)}{k}\right)^{(p-1)/p}\right],
 \quad
 \frac{|F_{\mathrm{ar},k}-F_{\sigma,k}|}{kq_k},
 \frac{|\Delta F_K|}{kq_k},\frac{|\Delta F_B|}{kq_k}
 =O_h\!\left[
 \left(\frac{\log(k+2)}{k}\right)^{(p-1)/p}\right].
 \tag{HMR44}
\]
To specify the rate substitution, in HMR38 take
$R=(k/\log(k+2))^{1/p}>1$. The original tilted tail bounds
$\log(M_\alpha/M_R)\leq
2C_h(1+R)^{1-p}/[(p-1)M_1]$. The explicit ceiling gives
$r=O_h(k)$ independently of $R$, while $Y\leq kR$.
For $N\in\{q_k-1,q_k,2q_k-1,2q_k\}$ one has
$N+1\geq(k+1)^2$ and $N=O_h(k^3)$.
Expanding the five nonnegative terms of
$\mathcal E_{N,r,Y}$ then gives
$\mathcal E_{N,r,Y}/[k(N+1)]
=O_h((1+R)\log(k+2)/k)$.
All other terms in HMR38 divided by $k(N+1)$ are bounded
by the displayed tilted-tail term plus
$O_h(\log(k+2)/k)$.
The two terms $R^{1-p}$ and $R\log(k+2)/k$ are both
exactly of the order in HMR44, and $(N_a+1)/q_k\leq2+1/q_k$.
This proves the claimed rate from the explicit finite formula;
the complete numerical constants are DMR30--33.
The finite HMR38 still holds at smaller degrees, but the
inequality $N+1\geq(k+1)^2$ is used here. In particular this
rate is not asserted at the simple-pair degree $N=2k+2$.